% !TEX encoding = UTF-8 Unicode

% This is a simple template for a LaTeX document using the "article" class.
% See "book", "report", "letter" for other types of document.
\documentclass[12pt,oneside]{article}
 % use larger type; default would be 10pt


\usepackage[utf8]{inputenc} % set input encoding (not needed with XeLaTeX)

%%% Examples of Article customizations
% These packages are optional, depending whether you want the features they provide.
% See the LaTeX Companion or other references for full information.

%%% PAGE DIMENSIONS
\usepackage{geometry}

 % to change the page dimensions
\geometry{a4paper} % or letterpaper (US) or a5paper or....
%\geometry{margins=2in} % for example, change the margins to 2 inches all round
%\geometry{landscape} % set up the page for landscape
%   read geometry.pdf for detailed page layout information

%\usepackage{graphicx} % support the \includegraphics command and options

% \usepackage[parfill]{parskip} % Activate to begin paragraphs with an empty line rather than an indent

%%% PACKAGES
\usepackage{booktabs} % for much better looking tables
\usepackage{array} % for better arrays (eg matrices) in maths
\usepackage{paralist} % very flexible & customisable lists (eg. enumerate/itemize, etc.)
\usepackage{verbatim} % adds environment for commenting out blocks of text & for better verbatim
\usepackage{subfig} % make it possible to include more than one captioned figure/table in a single float
% These packages are all incorporated in the memoir class to one degree or another...

%%% HEADERS & FOOTERS
%\usepackage{fancyhdr} % This should be set AFTER setting up the page geometry



%%% SECTION TITLE APPEARANCE

 % (See the fntguide.pdf for font help)
% (This matches ConTeXt defaults)

%%% ToC (table of contents) APPEARANCE
%\usepackage[nottoc,notlof,notlot]{tocbibind} % Put the bibliography in the ToC
%\usepackage[titles,subfigure]{tocloft} % Alter the style of the Table of Contents
\usepackage[encapsulated]{CJK} 


\usepackage{setspace}
\usepackage{amsfonts}
\usepackage{amssymb}
\usepackage{amsmath}
\usepackage{amstext}
\usepackage{amscd}
\usepackage{hyperref}
\usepackage{amsthm}

\newcommand{\e}[0]{\varepsilon}
\newcommand{\norm}[1]{\left|\left|#1\right|\right|}
\newcommand{\ra}[0]{\rightarrow}
\newcommand{\Ra}[0]{\Rightarrow}
\newcommand{\R}[0]{\mathbb{R}}
\newcommand{\Q}[0]{\mathbb{Q}}
\newcommand{\N}[0]{\mathbb{N}}
\newcommand{\D}[0]{\mathbb{D}}
\newcommand{\BF}[0]{\mathbb{F}}
\newcommand{\CN}[0]{\mathcal{N}}
\newcommand{\Z}[0]{\mathbb{Z}}
\newcommand{\C}[0]{\mathbb{C}}
\newcommand{\CE}[0]{\mathcal{E}}
\newcommand{\CM}[0]{\mathcal{M}}
\newcommand{\CA}[0]{\mathcal{A}}
\newcommand{\CP}[0]{\mathcal{P}}
\newcommand{\CT}[0]{\mathcal{T}}
\newcommand{\cond}[4]{\left \{ \begin{array}{ll} #1 & \mbox{#2} \\ #3 & \mbox{#4} \end{array} \right.}
\newcommand{\condd}[6]{\left \{ \begin{array}{ll} #1 & \mbox{#2} \\ #3 & \mbox{#4} \\ #5 & \mbox{#6}\end{array} \right.}
\newcommand{\conddd}[8]{\left \{ \begin{array}{ll} #1 & \mbox{#2} \\ #3 & \mbox{#4} \\ #5 & \mbox{#6} \\ #7 & \mbox{#8} \end{array} \right.}
\newcommand{\lams}[1]{\lambda^*(#1)}
\newcommand{\bs}[0]{\backslash}
\newcommand{\bM}[0]{\overline{\M}}
\newcommand{\bmu}[0]{\overline{\mu}}
\newcommand{\mus}[0]{\mu^*}
\newcommand{\sa}[0]{\sigma \mbox{-algebra}}
\newcommand{\set}[1]{\{ #1\}}
\newcommand{\CB}[0]{\mathcal{B}}
\newcommand{\CG}[0]{\mathcal{G}}
\newcommand{\CF}[0]{\mathcal{F}}
\newcommand{\CU}[0]{\mathcal{U}}
\newcommand{\CL}[0]{\mathcal{L}}
\newcommand{\FA}[0]{\mathfrak{A}}
\newcommand{\BT}[0]{\mathbb{T}}
\newcommand{\BI}[0]{\mathbb{I}}
\newcommand{\BE}[0]{\mathbb{E}}
\newcommand{\eK}[1]{\e\mbox{-Kronecker}(#1)}
\newcommand{\eKG}[0]{\e\mbox{-Kronecker}}
\newcommand{\KL}[2]{\mbox{{\renewcommand*\rmdefault{}\normalfont KL}}(#1 \ || \ #2)}
\newcommand{\JS}[2]{\mbox{{\renewcommand*\rmdefault{}\normalfont JS}}(#1 \ || \ #2)}

\newcommand{\IN}[1]{I_0(#1)}
\newcommand{\ING}[0]{I_0}
\newcommand{\st}[0]{such that }
\newcommand{\CS}[0]{\mathcal{S}}
\newcommand{\CC}[0]{\mathcal{C}}
\newcommand{\tn}[1]{\text{\normalfont #1}}



\newcommand{\bmp}[1]{\overline{M}^+(#1)}
\newcommand{\simplep}[2]{S_{#1}^+(#2)}
\newcommand{\simple}[2]{S_{#1}(#2)}
\newcommand{\sumf}[2]{\sum_{#1 = #2}^\infty}
\newcommand{\bigcapf}[2]{\bigcap_{#1 = #2}^\infty}
\newcommand{\bigcupf}[2]{\bigcup_{#1 = #2}^\infty}
\newcommand{\limf}[1]{\lim_{#1 \ra \infty}}
\newcommand{\seq}[2]{\set{#1_{#2}}_{#2=1}^\infty}
\newcommand{\sg}[1]{\sigma\langle#1\rangle}
\newcommand{\ag}[1]{\langle#1\rangle}
\newcommand{\g}[1]{\langle#1\rangle}
\newcommand{\inte}[0]{\mbox{int}}
\newcommand{\Cl}[0]{\mbox{cl}}
\newcommand{\bd}[0]{\mbox{bd}}
\newcommand{\inv}[1]{{#1}^{-1}}
\newcommand{\Aut}[0]{\mbox{{\renewcommand*\rmdefault{}\normalfont Aut}}}
\newcommand{\Gal}[0]{\mbox{{\renewcommand*\rmdefault{}\normalfont Gal}}}
\newcommand{\Bor}[0]{\mbox{{\renewcommand*\rmdefault{}\normalfont Bor}}}
\newcommand{\Meas}[0]{\mbox{{\renewcommand*\rmdefault{}\normalfont Meas}}}
\newcommand{\ball}[2]{\mbox{{\renewcommand*\rmdefault{}\normalfont B}}_{#1}({#2})}
\newcommand{\supp}[0]{\mbox{{\renewcommand*\rmdefault{}\normalfont supp}}}
\newcommand{\Trig}[0]{\mbox{{\renewcommand*\rmdefault{}\normalfont Trig}}}
\newcommand{\sign}[0]{\mbox{{\renewcommand*\rmdefault{}\normalfont sign}}}
\newcommand{\AP}[0]{\mbox{{\renewcommand*\rmdefault{}\normalfont AP}}}

\newcommand{\ali}[1]{\begin{align*}#1\end{align*}}
\newcommand{\zhongwen}[1]{\begin{CJK}{UTF8}{gbsn}#1\end{CJK}}
\renewcommand{\Bar}[1]{\overline{#1}}
\renewcommand{\Hat}[1]{\widehat{#1}}
\newcommand{\Ga}[0]{\Gamma}
\newcommand{\ga}[0]{\gamma}
\newcommand{\inner}[2]{\left\langle #1, #2 \right\rangle}
\renewcommand{\Tilde}[1]{\widetilde{#1}}


\theoremstyle{plain}
\newtheorem{thm}{Theorem}[subsection]
\newtheorem{lem}[thm]{Lemma}
\newtheorem{prop}[thm]{Proposition}
\newtheorem{cor}[thm]{Corollary}
\newtheorem{exmp}[thm]{Example}
\newtheorem*{defn}{Definition}



\theoremstyle{remark}
\newtheorem*{rem}{Remark}
\newtheorem*{claim}{Claim}
\newtheorem*{note}{Note}
\newtheorem*{sol}{Solution}
\newtheorem*{case}{Case}


%%% END Article customizations

%%% The "real" document content comes below...

\title{Background}
\date{} % Activate to display a given date or no date (if empty),
         % otherwise the current date is printed 
\parindent=0pt 
\begin{document}
\onehalfspace

\sloppy
\maketitle
\subsection*{Random Variables}
A real-valued {\it Random Variable} is a measurable function $X: E \ra \R$ from a probability space $(E, \CA, \nu)$ to the measure space $(\R, \CB, \lambda)$. 

Let $\mu_X$ be the push-forward $X^*\nu = \nu X^{-1}$. We call $\mu_X$ the distribution of $X$, denoted by $X \sim \mu_X$. This means for all $B \in \CB$, 
\ali{\nu(X^{-1}(B)) = \mu_X(B).} 

By the Law of Unconscious Statistician, 
\ali{\BE_{X \sim \mu_X}[X] = \int_\R x \ d\mu_X(x).}

More generally, if $f: \R \ra \R$ is a measurable function, then 
\ali{\BE_{X \sim \mu_X}[f(X)] = \int_\R f(x) \ d\mu_X(x).}

In case $\mu_X$ is absolute continuous with respect to $\lambda$, we let 
\ali{p_X = \frac{d\mu_X}{d\lambda} \in L^1(\R, \lambda)}
be the Radon-Nikodym derivative, and denote by $X \sim p_X$ with 
\ali{\BE_{X \sim p_X} [f(X)] = \int_\R f(x)p_X(x) \ d\lambda(x)}
for $f: \R \ra \R$ measurable. $p_X$ is the density function.

If $X$ is a random variable with values in a measure space, that is, $X: E \ra \Omega$ is a measurable function from a probability space $(E, \CA, \nu)$ to a measure space $(\Omega, \CM, \mu)$. Let $\mu_X = X^*\nu$ be the push-forward probability measure on $\Omega$, which is denoted by $X \sim \mu_X$.

$X$ itself does not have an expected value in general. If $f: \Omega \ra (\R, \CB, \lambda)$ is measurable, then 
\ali{\BE_{X \sim \mu_X}[f(X)] = \int_{\Omega} f(x) \ d\mu_X(x).}
Moreover, if $\mu_X$ is absolute continuous with respect to $\mu$, then we let 
\ali{p_X = \frac{d\mu_X}{d\mu} \in L^1(\Omega, \mu)}
be the Radon-Nikodym derivative, and we have, for $f: \Omega \ra (\R, \CB, \lambda)$ measurable,
\ali{\BE_{X \sim p_X} [f(X)] = \int_\Omega f(x)p_X(x) \ d\mu(x).}

\subsection*{Entropy}
\begin{defn}
Let $(\Omega, \CM)$ be a measurable space and $P$ and $Q$ be two probability measures on $\Omega$. Suppose $P$ is absolute continuous with respect to $Q$ (ie. $P \ll Q$). The Kullback-Leibler divergence (or K-L divergence, relative entropy), is defined as 
\ali{\KL{P}{Q} = \int_\Omega \log\left(\frac{dP}{dQ}(x)\right) \ dP(x),}
where $\frac{dP}{dQ} \in L^1(\Omega, Q)$ is the Radon-Nikodym derivative. 

More generally, suppose $P$ and $Q$ are both absolute continuous with respect to another measure $\mu$ (ie. $P \ll Q$ may not hold), we let $p = \frac{dP}{d\mu}$ and $q = \frac{dQ}{d\mu}$, both of which are in $L^1(\mu)$, and define
\ali{\KL{P}{Q} := \int_\Omega \log\left(\frac{p(x)}{q(x)}\right) p(x) \ d\mu(x).}
Note that such a measure $\mu$ always exists since we can put $\mu = \frac{P+Q}{2}$. 
\end{defn}

Note that $\KL{P}{Q} \geq 0$ because of Gibbs' inequality: Suppose $\set{p_1, p_2,..., p_n}$ and $\set{q_1, q_2,..., q_n}$ are discrete probability densities. Then
\ali{\sum_i p_i \log q_i \leq \sum_i p_i \log p_i,}
with equality if and only if $p_i = q_i$. Hence, $\KL{P}{Q} = 0$ if and only if $P = Q$ as measures or $p(x) = q(x)$ $\mu$-a.e.

\begin{defn}
Let $P$ and $Q$ be as above. Let $M = \frac{P+Q}{2}$. The Jensen-Shannon divergence is defined as 
\ali{\JS{P}{Q} := \frac{1}{2} \ \KL{P}{M} + \frac{1}{2} \ \KL{Q}{M}.}
\end{defn}

Note that $\JS{P}{Q} \geq 0$. It is known that $\JS{P}{Q} \leq \log 2$.


\end{document}